\documentclass[UTF8,11pt]{ctexart}
\usepackage[a4paper,margin=24mm]{geometry}
\usepackage{amsmath,amssymb,amsthm,booktabs,hyperref}
\newtheorem{proposition}{命题}
\newcommand{\ind}{\mathbf{1}}
\title{SPORK 下一轮完整 Actor 预计算：实现边界、正确性条件与收益推导}
\author{研究工作稿}
\date{2026年10月26日}
\begin{document}
\maketitle

本文区分代码事实、条件性推导、既有文献和待验证假设。推导用于解释当前实现，
不据此宣称新的投机执行定理或顶会级原创性。参考基线为相同服务配置下的 SPORK D1/D2/D3；
当前方法名为 raw。本稿不修改已冻结代码、实验记录或评测协议。

\section{当前实现与未采用的分支}
当前方法保留父轮完整推理。SPORK 的部分思考前缀探针产生候选工具调用，
工具在独立环境中真实执行。候选结果先于父轮推理结束时，构造候选下一轮输入，
在同一目标模型上提前运行完整的下一轮 Actor。它的工具调用也在独立环境中执行。
父轮正常结束后，复用条件全部成立才接纳下一轮执行；否则取消、回收并执行正式下一轮。

真实观测指“执行候选动作实际得到的观测”，不表示动作已获父轮授权。
输入来自已提交历史、当前可见前缀产生的候选动作及实际结果、环境副本和模板，
没有预读正确未来轨迹。输出是可条件复用的下一轮完整响应及执行后状态。
没有新训练、训练标签或在线参数更新。

当前只保留一个候选下一轮，不递归展开；按 $(1,2),(3,4),\ldots$ 配对。
可同时存在两条完整 Actor 流和原 SPORK 的短探针。
父轮保留 D3，未提交子轮保留 D1/D2、关闭 D3，因为现有清理接口作用于整个服务。
这一点是实现约束，不能作为新的 token-level speculative decoding 算法。
最终提交前再次检查父状态和输入；取消时等待工具线程收尾，避免旧分支后续写回。

目录中保留了其他分支，但它们没有用于本次 310 次全量实验：
32/96-token BERT 启动控制器、ETA 式提前结束思考、候选刷新与影子回滚、
额外确认探针 confirmed。confirmed 只做过小样本比较，不足以否定所有确认策略。
经验 memory/pattern 检索和隐藏状态分类器也未接入当前方法。

\section{下一轮输入的依赖关系}
设 $H_t$ 为已提交消息历史，$S_t$ 为环境完整相关状态，
$X_t=\operatorname{Encode}(\operatorname{Render}(H_t))$ 为实际模型输入。
$g_t$ 为生成配置，$\xi_t$ 为相关随机性。一轮正式转移为
\begin{align}
Y_t &= F_\theta(X_t;g_t,\xi_t), & a_t &= \operatorname{Parse}(Y_t),\\
(o_t,S_{t+1}) &= E(S_t,a_t), & H_{t+1} &= U(H_t,Y_t,o_t).
\end{align}
终止回答单独结束循环，不假设存在下一条工具依赖边。
$U$ 接收整个正式输出 $Y_t$，而不一定仅接收动作 $a_t$。
本地基线去掉历史消息开头的思考块，但工具调用前的可见文字、格式残留等仍可影响渲染。
因此不能一般地断言“动作和观测相同，所以下一轮输入相同”。

投机动作 $\hat a_t$ 从当前状态副本执行：
\[
(\hat o_t,\hat S_{t+1})=E(\operatorname{Clone}(S_t),\hat a_t).
\]
以候选工具消息 $\hat Y_t$ 构造 $\hat H_{t+1}=U(H_t,\hat Y_t,\hat o_t)$ 和 $\hat X_{t+1}$。
这些量在启动时来自已可见信息；未完成的正式输出仅在随后验证时使用。

\begin{proposition}[输入可提前构造的条件]
若历史更新和模板在当前分支上只依赖已知历史、已预测动作和已获得观测，
不依赖父轮尚未生成的内容，则该分支的下一轮输入可以提前构造。
若未知内容仍会被保留到下一轮输入中，动作相等本身不足以授权复用。
\end{proposition}
\begin{proof}
前者是对已知变量计算确定函数。后者可取动作相同、保留文字不同的两个父轮输出，
它们经 $U$ 和渲染可以产生不同输入，故动作相同不能推出输入相同。
\end{proof}
这是依赖关系的形式化，不是既有投机执行理论的原创性主张。
不要求所有候选都满足条件：未满足者由实际完整输入检查拒绝。

\section{复用条件与条件性正确性}
当前验证条件为
\[
I_t=\ind\!\left[
\hat a_t=a_t\land\hat o_t=o_t\land\hat S_{t+1}\equiv S_{t+1}
\land\hat X_{t+1}=X_{t+1}\land\hat g_{t+1}=g_{t+1}
\right].
\]
动作按工具名及原 SPORK 参数规则比较，观测使用 SHA-256，输入比较完整 token ID 列表。
代码的状态键仅含 agent/user 数据库摘要；生成键含模型标识、地址、seed、预算、
上下文长度、thinking、stop 和 temperature。
这不是对任意外部状态、权重变化、LoRA、logits processor 和随机源的完整覆盖证明。

\begin{proposition}[相对于参考执行的轨迹一致性]
假设：
(i) 固定输入和配置的模型输出与调度无关，或可与参考逐请求耦合同一随机流；
(ii) 环境转移由所验证状态、动作及所耦合随机性完全决定；
(iii) 所有投机副作用均隔离并且只提交一次；
(iv) 状态指纹等价于相关状态等价，忽略摘要碰撞；
(v) 验证至提交之间无未检查变化，失败和截断得到一致处理；
(vi) 原 SPORK D3 保持参考模型语义。
则接纳 $I_t=1$ 的下一轮，或在 $I_t=0$ 时执行正式下一轮，保持参考的已提交轨迹。
\end{proposition}
\begin{proof}
对已提交轮次归纳。初始消息和环境相同。父轮完整运行，在假设下决定相同。
拒绝时，隔离分支无外部效应，下一轮从真实历史和状态执行。
接纳时，完整输入、配置及起始相关状态相同，子轮模型输出和工具转移因而相同；
其状态只提交一次，且等待结束后再次检查父状态和输入。
两种路径都建立下一归纳步。
\end{proof}
greedy/batch-invariant 下的实测 token 一致性支持所测范围，不能替代任意后端的证明。
同 seed 不保证调度不变，数据库 hash 不覆盖任意外部状态，deepcopy 不能隔离任意外部支付或网络效应。
随机模型的分布性扩展还要求不按子轮答案内容或完成速度筛选随机输出。
正确性依赖验证与隔离，不依赖预测器总是正确。

\section{理想增量收益}
从一对相邻轮次开始计时。$R$ 为父轮完整推理/动作生成时间，$d$ 为候选工具发起时间，
$L$ 为父轮工具执行时间，$B$ 为参考 SPORK 下一轮的完整耗时。
先假设候选正确、$0\le d\le R$、无竞争及管理开销，下一轮工作量不变。
\begin{align}
T_{\mathrm{SPORK}} &= \max(R,d+L)+B,\\
T_{\mathrm{next}} &= \max(R,d+L+B).
\end{align}
观测在 $d+L$ 到达，父轮在 $R$ 结束，因此新增窗口为 $W=(R-d-L)_+$。
分情况相减得到
\[
\boxed{\Delta T_{\mathrm{ideal}}=\min(B,W)=\min\{B,(R-d-L)_+\}.}
\]
若 $d+L\ge R$，无父轮推理窗口；否则最多隐藏完整下一轮的 $B$ 时间。
同一理想模型下，原 SPORK 相对于串行父轮最多隐藏 $\min\{L,R-d\}$；
扩展后父轮工具和下一轮合计隐藏 $\min\{L+B,R-d\}$。
新增收益不受父轮工具时长 $L$ 单独限制。

例如 $R=100,d=18,L=2,B=80$ 秒，SPORK 为 $180$ 秒，新方案为 $100$ 秒，
理想增量为 $80$ 秒。这是解释依赖关系的例子，不是实测，也未计入失败与 GPU 竞争。
其他量固定时，较短父轮工具时间可扩大新模块窗口；不能推广成所有模型、工具和 GPU 条件下的单调规律。
对于独立两轮组，理想新耗时不小于 $\max(A,B)$，$A=\max(R,d+L)$，
所以相对于 $A+B$ 的加速不超过 $2$。
此界依赖固定工作量，不适用于任意 cache/D3 变化及运行噪声条件。
类似最小值重叠公式及深度界已有 SA、CTES、ProSpect 等覆盖，公式本身不构成新的理论贡献。

\section{共享 GPU、拒绝与取消成本}
参考两轮耗时为 $A+B$。接受时，实际父轮耗时为 $A'$，
子轮在 $r\le A'$ 发起，实际持续时间为 $B'$。
开始等待子轮前的验证耗时为 $v_1$，提交及剩余回收时间为 $v_2$：
\[
T_{\mathrm{hit}}=\max(A'+v_1,r+B')+v_2.
\]
令 $O=\min\{B',(A'-r)_+\}$，
$e=\max(A'+v_1,r+B')-\max(A',r+B')+v_2$，有 $v_2\le e\le v_1+v_2$，因此
\[
G_{\mathrm{hit}}=A+B-T_{\mathrm{hit}}
=O-(A'-A)-(B'-B)-e.
\]
$A'-A$ 包含竞争和处理造成的父轮变化；$B'-B$ 还含子轮关闭 D3 的影响。
cache 与调度可能使某些差值为负，不能将每个差值预设为正的竞争损失。
拒绝时，设回收耗时为 $K$、回收后正式下一轮耗时为 $B_f$，其余开销为 $V_m$。
当前代码先回收再继续：
\[
T_{\mathrm{miss}}=A'+K+B_f+V_m,\qquad
G_{\mathrm{miss}}=-(A'-A)-(B_f-B)-K-V_m.
\]
父轮直接回答时令 $B=B_f=0$，仍计入错误子轮的竞争和回收。
用 $I$ 表示接受事件，$V=Ie+(1-I)V_m$，则
\[
\boxed{\mathbb E[G]=\mathbb E[I\,O]-
\mathbb E[\Delta A+I\Delta B+(1-I)\Delta B_f+V+(1-I)K].}
\]
必须保留联合项 $\mathbb E[IO]=P(I=1)\mathbb E[O\mid I=1]$；
窗口与接受概率一般相关，不能无条件写成 $P(I=1)\mathbb E[O]$。
不同运行的缓存和测量噪声也限制了分项的因果识别。
将平均额外开销记为 $\bar C$，净加速条件是 $p\mathbb E[O\mid I=1]>\bar C$。
若接受时净收益至少为 $g_{\min}>0$、拒绝惩罚至多为 $c_{\max}\ge0$，
则 $p>c_{\max}/(g_{\min}+c_{\max})$ 是一个充分条件。
该条件需要可信概率和成本界；模型置信度不自动构成这种界。

\section{算力成本与可选控制}
接纳的子轮替代原本要执行的正式下一轮，不能把它全部算成新增计算。
概念上：
\[
\Delta C=\Delta C_{\mathrm{parent}}
+I(C_{\mathrm{child}}-C_{\mathrm{next}})
+(1-I)(C_{\mathrm{waste}}+C_{\mathrm{fallback}}-C_{\mathrm{next}})
+C_{\mathrm{admin}}.
\]
额外成本包含错分支、prefill/探针、D3 差异和管理，也可能被 cache 复用部分抵消。
记录的被丢弃 token 只是已接收的错误分支 decode 下界，
不覆盖所有取消中的 GPU 工作、prefill、能耗及 GPU 忙碌时间。

未来若控制启动与否，特征 $x$ 只能取当前已知历史、前缀、候选调用、已返回观测、
已用时间、估计剩余生成量、队列与显存等。概念目标为
\[
\mathbb E[G\mid x]-\lambda\mathbb E[\Delta C\mid x]>0.
\]
$\lambda$ 将算力成本换算到延迟目标。期望效用规则已有相关工作，
不能作为控制器已经实现或有效的证据。
接受概率标签可由未来验证结果给出；净收益标签需要同条件 on/off 对照。
未来时长不能进入运行时输入；失败和截断也不能简单当作普通负样本。
memory 只能读取此前已完成任务，开发、验证和最终评测必须区分。

当前不宜优先重启 BERT 控制器：104 个正常训练任务的 289 个两轮段上，
事后逐段选择更快方案比始终使用当前方法额外省 $567.15$ 秒，
即 $5.45$ 秒/任务，约 $1.82\%$。这只是乐观经验参照；
改变段选择会改变后续 cache 和调度，不能视为因果策略收益或严格上界。
应先验证核心机制和算力代价，再决定是否训练预算控制。

\section{已有证据及边界}
正式实验为 155 任务、两配置、每配置一次，共 310 次运行。
148 个共同无运行错误的任务平均耗时 $364.004\to293.334$ 秒，减少 $19.415\%$；
22 个有效测试任务减少 $17.688\%$。
148 对正式 Actor token、动作、工具结果和记录状态一致。
两组在相同 7 个任务触及长度上限，12 个有效配对触及 8 轮上限；
去掉后者，136 对任务仍减少 $19.087\%$。
固定用户消息及 ACTION reward 不等于官方完整交互成功率。

有效任务中发起 405 个子轮、接纳 263 个。
这些接受分支从发起事件到父轮结束平均有 $62.52$ 秒窗口。
发起事件在 asyncio 任务创建前，不是 GPU 首 token 时间，也不是反事实收益上界。
219 个接受子轮含工具调用，18 个改变了记录数据库状态。
此结果支持所测环境的隔离提交，不能扩展到未覆盖的外部效应。

7 次分支通过动作、观测、父状态检查，却因完整下一轮输入不同而被拒绝：
2 次有额外可见说明，5 次有错误工具结束标签残留。
后 5 次属于格式/实现边界，不能包装成研究贡献；
前 2 次说明正常可见文字也能破坏“同动作推出同输入”。
检查全部保留；这些数据不证明删除其他检查安全。

\section{贡献定位}
SPORK~\cite{spork} 已有前缀探针、提前工具执行及候选 token 的 D3。
SA~\cite{sa} 已把 LLM/API 调用统一为可投机操作；
ProSpect 本地稿件也明确提前执行下一次目标模型调用并按输入匹配接纳。
ObsSpec~\cite{obsspec} 预测观测并提前运行 policy；
AOSpec~\cite{aospec} 结合观测投机、隔离动作执行和动作/起始状态验证。
CTES 本地稿件已有受控多步展开、兼容状态复用及重叠/成本分析。

可争取的主张集中为：长思考 Agent 中真实候选观测已经到达、父轮仍未结束的执行机会；
完整下一轮输入可构造、可验证的条件；同一在线 GPU 上保持已提交轨迹的跨轮完整目标模型执行；
以及收益随竞争、工具时长和算力成本变化的适用区域。
这些是需要强对照和跨环境证据的贡献候选，不是已确认首次性。
若同预测输入和资源配置的 SA/ProSpect 式实现与当前方法实质等价，
应承认方法层面主要是 SPORK 适配，不能用命名或模块拆分制造新颖性。

\section{实验优先级}
\begin{enumerate}
\item 机制消融：训练划分每领域最小四个任务 ID，共 8 任务，最多两轮、两次交错顺序重复。
四组为原 SPORK、串行且子轮关闭 D3 的等配置对照、只提前子轮模型输出而延后工具执行、
当前完整方法，共 64 次运行。新消融需要先实现及 CPU 检查，当前未运行。
\item 适用区域：逐因素改变工具延迟下限、背景负载和思考长度，测窗口、父轮变慢、
prefill/decode、取消及显存。全局 D3 清理限制带 D3 的多 Agent 并发；
背景负载先采用不带 D3 的请求，或先实现按请求隔离。
\item 重复性：冻结实现后追加独立重复并交换顺序，保留所有失败。
跨任务 bootstrap 不代替运行噪声评估。若使用已看过的测试结果调参，
后续需新的未使用任务作为最终评测。
\item 完整交互：先核查用户模拟器协议与回复，再用完整用户交互及官方评估器；
分开报告 Actor、用户模型和总耗时，对照保持资源一致。
\item 泛化及强基线：补另一目标模型和另一任务环境；本地 Qwen3-4B 是候选，
serving 配置仍需核查。对齐预测输入、真实工具、上下文规则、GPU 预算及计时后比较近邻方法。
\item 可选扩展：先证实成本或饱和负载是限制，再做预算控制、memory 概率先验，
或按请求隔离的子轮 D3，避免同时堆叠模块。
\end{enumerate}
所有 GPU 和耗时实验由用户在 tmux 中启动。
本稿仅进行文献核对、已有记录的轻量统计和理论整理。

\begin{thebibliography}{9}
\bibitem{spork} Bai et al. SPORK: Self-Speculative Forking to Accelerate Agentic LLM Inference.
\url{https://arxiv.org/abs/2607.03333}.
\bibitem{sa} Ye et al. Speculative Actions: A Lossless Framework for Faster Agentic Systems.
\url{https://arxiv.org/abs/2510.04371}.
\bibitem{obsspec} Montgomery et al. ObsSpec: Accelerating Agentic Reinforcement Learning with Observation Speculation.
\url{https://arxiv.org/abs/2610.15246}.
\bibitem{aospec} Chen et al. AOSpec: Action and Observation Co-Speculation for Low-Latency Agent Serving.
\url{https://arxiv.org/abs/2608.00881}.
\end{thebibliography}
\end{document}
